\section{Benchmark Construction Details}
\label{app:construction}

\subsection{Graph-Authoring Workflow}

The graph-authoring workflow has four reviewed stages:
\begin{enumerate}[leftmargin=1.7em]
    \item \textbf{Node extraction.} Segment-level calls propose diagnosable units with source locators and grounding notes.
    \item \textbf{Reconciliation.} Candidate nodes are merged, split, and deduplicated while preserving provenance.
    \item \textbf{Diagnostic enrichment.} Each node receives a definition, diagnostic goal, rubric, diagnostic signals, and simulator guidance.
    \item \textbf{Edge proposal and review.} Relations are proposed from the complete node set and manually reviewed before promotion.
\end{enumerate}

The relation vocabulary includes \texttt{part\_of}, \texttt{prerequisite\_for}, \texttt{supports}, and \texttt{contrasts\_with}. Candidate confidence does not replace human review.

\subsection{Knowledge-Map Review}

Reviewers verify that each candidate map covers every node and that each state is supported by behaviorally realizable evidence.

They also check consistency among mastery, misconceptions, unknowns, profile context, and graph structure. Promotion creates an immutable reviewed snapshot.

\subsection{Manifest-Bound Runtime Contexts}

The runtime maintains three contexts:
\begin{itemize}[leftmargin=1.5em]
    \item \textbf{Tested-agent context:} public graph, interaction rules, visible transcript, and working map.
    \item \textbf{Simulator context:} grounded hidden states, evidence, limited visible dialogue, and optional style context.
    \item \textbf{Evaluation context:} hidden map, final reconstruction, and deterministic scoring configuration.
\end{itemize}

An episode manifest fixes the artifact versions and execution settings shared by these contexts. Debug traces and hidden policy outputs never enter the tested-agent context.
